%============================================================
% CH1 · P5 · PROBLEM 3 — FIND THE DUPLICATE NUMBER (LeetCode 287)
%============================================================
\subsection{Problem 1.5.3: Find the Duplicate Number}
\problemlabel{1.5.3}{ssec:ch1-find-duplicate}

\problemcard{1.5.3}{Find the Duplicate Number}{LeetCode 287}{https://leetcode.com/problems/find-the-duplicate-number/}{\exmedium\ \lc{Medium}}{Read the array as arrows $i \to a[i]$; the walk from $0$ is a $\rho$ whose entry point is the duplicate; Floyd's tortoise and hare finds it with \bigO{1} memory}

\textbf{What this problem adds.} Problems~1.5.1 and~1.5.2 marked vertices in a \texttt{visited} or \texttt{when} array. Here that is forbidden: the array may not be modified and only \bigO{1} extra memory is allowed, so Algorithm~4 of the theory (Floyd's tortoise and hare) is no longer an optional curiosity but the whole point. The graph is hidden in an array of numbers, and the answer is the \textbf{entry point} of the $\rho$ that starts at vertex $0$ --- this time we also prove carefully why the second phase of Floyd's method lands exactly there.

\subsubsection*{1. Problem statement}

An array \texttt{nums} has $n+1$ integers, each in the range $[1, n]$. Exactly one value occurs more than once (it may occur two or more times); return that value. You must \textbf{not modify} \texttt{nums} and may use only \textbf{constant extra space}. Follow-up questions: prove that a repeated value must exist, and solve the problem in less than \bigO{n^2} time.

\textbf{Constraints:} $1 \le n \le 10^5$, \; $\texttt{nums.length} = n+1$, \; $1 \le \texttt{nums}[i] \le n$, \; all values except one appear at most once.

\begin{table}[H]
\centering\small
\begin{tabular}{@{}lcl@{}}
\toprule
\textbf{nums} & \textbf{Output} & \textbf{Explanation} \\
\midrule
$[1,3,4,2,2]$ & 2 & $2$ appears at indices $3$ and $4$ \\
$[3,1,3,4,2]$ & 3 & $3$ appears at indices $0$ and $2$ \\
$[3,3,3,3,3]$ & 3 & one value five times; $1, 2, 4$ are missing \\
\bottomrule
\end{tabular}
\end{table}

\begin{diagrambox}
\begin{tikzpicture}
    \begin{scope}
        \node[font=\small\bfseries] at (1.7,0.9) {Example 1: $[1,3,4,2,2]$};
        \node[gvsrc] (a0) at (0,0) {0};
        \node[gv] (a1) at (1.1,0) {1};
        \node[gv] (a3) at (2.2,0) {3};
        \node[gvdst] (a2) at (3.3,0) {2};
        \node[gvseen] (a4) at (3.3,-1.3) {4};
        \draw[gde] (a0) -- (a1); \draw[gde] (a1) -- (a3); \draw[gde] (a3) -- (a2);
        \draw[gde] (a2) to[bend left=35] (a4); \draw[gde] (a4) to[bend left=35] (a2);
    \end{scope}
    \begin{scope}[xshift=5.1cm]
        \node[font=\small\bfseries] at (1.2,0.9) {Example 2: $[3,1,3,4,2]$};
        \node[gvsrc] (b0) at (0,0) {0};
        \node[gvdst] (b3) at (1.2,0) {3};
        \node[gvseen] (b4) at (2.5,0) {4};
        \node[gvseen] (b2) at (1.85,-1.1) {2};
        \node[gvdead] (b1) at (0,-1.1) {1};
        \draw[gde] (b0) -- (b3); \draw[gde] (b3) -- (b4); \draw[gde] (b4) -- (b2); \draw[gde] (b2) -- (b3);
        \draw[gde] (b1) to[out=200, in=250, looseness=6] (b1);
    \end{scope}
    \begin{scope}[xshift=9.9cm]
        \node[font=\small\bfseries] at (1.2,0.9) {Example 3: $[3,3,3,3,3]$};
        \node[gvsrc] (c0) at (0,0) {0};
        \node[gv] (c1) at (0,-1.1) {1};
        \node[gv] (c2) at (2.4,0) {2};
        \node[gv] (c4) at (2.4,-1.1) {4};
        \node[gvdst] (c3) at (1.2,-0.55) {3};
        \foreach \v in {0,1,2,4} \draw[gde] (c\v) -- (c3);
        \draw[gde] (c3) to[out=240, in=300, looseness=6] (c3);
    \end{scope}
\end{tikzpicture}
\figcap{Each index $i$ is a vertex with one arrow to \texttt{nums[i]}. Green: vertex $0$, which no arrow can enter (all values are $\ge 1$). Orange: the duplicate --- in every example it is the vertex where the walk from $0$ enters its cycle. In Example 2 the self-loop at $1$ (grey) is a separate piece that the walk from $0$ never sees.}
\end{diagrambox}

\subsubsection*{2. The graph model}

\begin{modelbox}[Graph Model --- Find the Duplicate Number]
\renewcommand{\arraystretch}{1.35}
\begin{tabularx}{\linewidth}{@{}>{\bfseries}l X@{}}
Vertices & the $n+1$ indices $0..n$ \\
Edges & $i \to \texttt{nums}[i]$; every value is a valid index ($1 \le \texttt{nums}[i] \le n$), so out-degree is exactly 1: a \textbf{functional graph}, stored in the read-only input array \\
Weights & none \\
Question & the entry point of the $\rho$ that starts at vertex $0$ (= the vertex with two arrows into it on that walk) \\
Size & $|V| = n+1 \le 10^5+1$, $|E| = |V|$; memory \bigO{1} besides the input, input read-only; time below \bigO{n^2} \\
\end{tabularx}
\end{modelbox}

\begin{diagrambox}
\begin{tikzpicture}
    \dpindex{ix}{0.75}{$i$}{0}{4}
    \dparray{va}{0}{$a[i]$}{1/cell, 3/cell, 4/cell, 2/anscell, 2/anscell}
    \begin{scope}[xshift=6.2cm]
        \node[gvsrc] (m0) at (0,0) {0};
        \node[gv] (m1) at (1.25,0) {1};
        \node[gv] (m3) at (2.5,0) {3};
        \node[gvdst] (m2) at (3.75,0) {2};
        \node[gvseen] (m4) at (5.25,0) {4};
        \draw[gde] (m0) -- (m1); \draw[gde] (m1) -- (m3);
        \draw[gpath, -{Stealth[length=3mm]}] (m3) -- node[lbl, below=4mm]{$a[3]{=}2$} (m2);
        \draw[gde] (m2) to[bend left=40] (m4);
        \draw[gpath, -{Stealth[length=3mm]}] (m4) to[bend left=40] node[lbl, below]{$a[4]{=}2$} (m2);
        \draw[decorate, decoration={brace, amplitude=4pt}, darkgray] (-0.3,0.55) -- (2.8,0.55)
            node[midway, above=4pt, font=\scriptsize]{tail $\mu = 3$};
        \draw[decorate, decoration={brace, amplitude=4pt}, darkgray] (3.45,0.75) -- (5.55,0.75)
            node[midway, above=4pt, font=\scriptsize]{cycle $\lambda = 2$};
    \end{scope}
\end{tikzpicture}
\figcap{Example 1 as an array and as a $\rho$. The entry point $2$ receives one arrow from the tail ($a[3] = 2$) and one from the cycle ($a[4] = 2$): two different indices hold the value $2$ --- that \emph{is} the duplicate.}
\end{diagrambox}

\subsubsection*{3. Key observations}

\begin{itemize}[leftmargin=*]
    \item \textbf{A duplicate exists (pigeonhole).} $n+1$ values in the $n$ boxes $1..n$: some box gets two. This answers the first follow-up.
    \item \textbf{Vertex $0$ is never on a cycle.} Every cycle vertex has an incoming arrow from its cycle predecessor, but no value equals $0$. So the walk from $0$ has a tail of length $\mu \ge 1$.
    \item \textbf{The entry point has two incoming arrows.} The entry $e$ is entered from the last tail vertex \emph{and} from the last cycle vertex; these are two different indices with value $e$. Since only one value repeats, $e$ is the answer.
    \item \textbf{The start matters.} From a vertex that lies on a cycle (index $1$ in Example 2) the walk has no tail, its first repeated vertex is the start itself, and that value need not be a duplicate. Vertex $0$ is the one start that is \emph{guaranteed} to have a tail.
    \item \textbf{The constraints forbid the earlier tools.} \texttt{first[]}, \texttt{state[]}, \texttt{when[]} all cost \bigO{n} memory, and marking inside \texttt{nums} (e.g.\ negating entries) modifies the input. Only a constant number of vertex variables may be kept: pointers that walk.
\end{itemize}

\begin{diagrambox}
\begin{tikzpicture}
    \begin{scope}
        \node[font=\small\bfseries] at (1.4,1.0) {walk from $0$: tail, then cycle};
        \node[gvsrc] (p0) at (0,0) {0};
        \node[gvdst] (p3) at (1.3,0) {3};
        \node[gvseen] (p4) at (2.7,0.55) {4};
        \node[gvseen] (p2) at (2.7,-0.55) {2};
        \draw[gpath, -{Stealth[length=3mm]}] (p0) -- (p3);
        \draw[gde] (p3) -- (p4); \draw[gde] (p4) -- (p2);
        \draw[gpath, -{Stealth[length=3mm]}] (p2) -- (p3);
        \node[note, text width=4.4cm] at (1.4,-1.4) {$3$ is entered from $0$ (tail) and from $2$ (cycle): $a[0] = a[2] = 3$, \textbf{a duplicate}};
    \end{scope}
    \begin{scope}[xshift=7.6cm]
        \node[font=\small\bfseries] at (0.6,1.0) {walk from $1$: no tail};
        \node[gvdead] (q1) at (0,0) {1};
        \draw[gde] (q1) to[out=30, in=-30, looseness=6] (q1);
        \node[note, text width=4.6cm] at (0.6,-1.4) {the first repeated vertex is the start $1$ itself, entered only by its own arrow: $1$ appears once, \textbf{not} a duplicate};
    \end{scope}
\end{tikzpicture}
\figcap{Example 2 seen from two starts. Only a start outside every cycle --- and $0$ is always one --- produces an entry point with two incoming arrows.}
\end{diagrambox}

\subsubsection*{4. Approaches}

\begin{bruteforcebox}[Approach 1: Compare All Pairs]
Respect both memory rules in the most naive way: compare every pair of positions.

\begin{lstlisting}[style=cpp]
#include <bits/stdc++.h>
using namespace std;

class Solution {
public:
    int findDuplicate(vector<int>& nums) {
        int m = nums.size();              // m = n + 1
        for (int i = 0; i < m; i++)
            for (int j = i + 1; j < m; j++)
                if (nums[i] == nums[j]) return nums[i];   // read-only, O(1) memory
        return -1;                        // unreachable: a duplicate always exists
    }
};
\end{lstlisting}

\begin{complexitybox}
\textbf{Time:} \bigO{n^2} --- about $5\cdot10^9$ comparisons for $n = 10^5$ when the duplicate pair sits at the end. \textbf{Space:} \bigO{1}, input untouched. \textbf{Verdict:} \tle{}; it obeys the memory rules but fails the time follow-up. It serves as the reference for random tests.
\end{complexitybox}
\end{bruteforcebox}

\begin{algorithmbox}[Approach 2: Floyd's Tortoise and Hare from Vertex 0]
Use template (5) of the theory with $s = 0$ and $f(i) = \texttt{nums}[i]$, but return the entry point instead of $\mu$ and $\lambda$.
\begin{enumerate}[leftmargin=*]
    \item \textbf{Phase 1:} tortoise $t$ moves 1 step, hare $h$ moves 2 steps per round, until $t = h$ (a vertex on the cycle).
    \item \textbf{Phase 2:} put $t$ back on $0$; both move 1 step per round; where they meet is the entry point $=$ the duplicate.
\end{enumerate}

\begin{lstlisting}[style=cpp]
#include <bits/stdc++.h>
using namespace std;

class Solution {
public:
    int findDuplicate(vector<int>& nums) {
        // functional graph i -> nums[i] on 0..n; no arrow enters 0
        int t = nums[0], h = nums[nums[0]];   // round 1: tortoise 1 step, hare 2 steps
        while (t != h) {                      // phase 1: meet somewhere on the cycle
            t = nums[t];
            h = nums[nums[h]];
        }
        t = 0;                                // phase 2: tortoise back to the start
        while (t != h) {                      // both 1 step: they meet at the entry
            t = nums[t];
            h = nums[h];
        }
        return t;                             // entry of the cycle = the duplicate value
    }
};
\end{lstlisting}

\begin{complexitybox}
\textbf{Time:} \bigO{n} --- phase 1 takes fewer than $\mu + \lambda \le n+1$ rounds, phase 2 exactly $\mu$ steps (proof below). \textbf{Space:} \bigO{1}: two integers; \texttt{nums} is only read. \textbf{Verdict:} \ac{}, and it meets every requirement of the statement.
\end{complexitybox}
\end{algorithmbox}

\begin{invariantbox}[Why it is correct --- the $\mu$, $\lambda$ argument]
Write $x_k = f^k(0)$ for the vertex after $k$ steps. By the shape theorem $x_0, \dots, x_{\mu+\lambda-1}$ are all different, $x_\mu = e$ is the entry point, and $x_{k+\lambda} = x_k$ for every $k \ge \mu$. A tail vertex ($k < \mu$) appears only once in the whole sequence.

\textbf{(a) The entry point is the duplicate.} $\mu \ge 1$ because $0$ is not on a cycle. The vertex $e$ is the successor of $x_{\mu-1}$ (a tail vertex) and of $x_{\mu+\lambda-1}$ (a cycle vertex; for $\lambda = 1$ it is $e$ itself). These indices differ, and both hold the value $e$. Exactly one value repeats, so it is $e$.

\textbf{(b) Phase 1 meets.} After $t$ rounds the tortoise is at $x_t$ and the hare at $x_{2t}$. For $t < \mu$, $x_t$ is a tail vertex that never appears again, so $x_{2t} \ne x_t$. For $t \ge \mu$ both are on the cycle, and $x_{2t} = x_t$ iff $\lambda$ divides $t$. So the first meeting happens after $t^\ast$ rounds, the smallest positive multiple of $\lambda$ with $t^\ast \ge \mu$; since $\mu \ge 1$, $t^\ast \le \mu + \lambda - 1 \le n$.

\textbf{(c) Phase 2 meets at $e$, and nowhere earlier.} After $k$ more steps the restarted pointer is at $x_k$ and the other one at $x_{t^\ast + k}$. For $k = \mu$: $x_{t^\ast + \mu} = x_\mu$, because $t^\ast$ is a multiple of $\lambda$ and whole laps after the tail change nothing. For $k < \mu$: $x_k$ is on the tail, but $x_{t^\ast + k}$ has index $\ge t^\ast \ge \mu$, so it is on the cycle; they differ. Hence the pointers first coincide after exactly $\mu$ steps, at $x_\mu = e$. \hfill$\blacksquare$
\end{invariantbox}

\begin{diagrambox}
\begin{tikzpicture}
    \node[gvsrc] (z0) at (0,0) {0};
    \node[gv] (z1) at (1.2,0) {};
    \node[gv] (z2) at (2.4,0) {};
    \draw[gde] (z0) -- (z1); \draw[gde, dashed] (z1) -- (z2); \draw[gde] (z2) -- (3.25,0);
    \draw[gpath, -{Stealth[length=3mm]}, shorten >=3.6mm] (3.6,0) arc[start angle=180, end angle=300, radius=1.3];
    \draw[softpurple, line width=1.4pt, -{Stealth[length=3mm]}, shorten >=3.6mm] ($(4.9,0)+(300:1.3)$) arc[start angle=300, end angle=540, radius=1.3];
    \node[gvdst] (ze) at (3.6,0) {$e$};
    \node[gvcur] (zm) at ($(4.9,0)+(300:1.3)$) {$m$};
    \node[lbl, text=accentorange] at ($(4.9,0)+(235:1.75)$) {$d$};
    \node[lbl, text=softpurple] at ($(4.9,0)+(60:1.75)$) {$\lambda - d$};
    \draw[decorate, decoration={brace, amplitude=5pt, mirror}, darkgray] (-0.3,-0.6) -- (3.3,-0.6)
        node[midway, below=5pt, font=\scriptsize]{tail: $\mu$ steps from $0$ to $e$};
    \node[note, text width=6.4cm, align=left, anchor=west] at (6.9,0.1)
        {Phase 1 ends at $m = x_{t^\ast}$, which is $d = t^\ast - \mu$ steps after $e$ on the cycle. Since $\lambda \mid t^\ast$, $\mu \equiv -d \equiv \lambda - d \pmod{\lambda}$: walking $\mu$ steps from $m$ means $\lambda - d$ steps back to $e$ plus whole laps --- and the pointer from $0$ needs the same $\mu$ steps to reach $e$.\\[2pt]
        Trace input of Section 5: $\mu = 4$, $\lambda = 3$, $t^\ast = 6$, $d = 2$, and $4 = (3-2) + 1\cdot 3$.};
\end{tikzpicture}
\figcap{The picture behind phase 2. Orange arc: the $d$ steps from $e$ to the meeting point $m$. Purple arc: the $\lambda - d$ steps from $m$ back to $e$; the remaining $\mu - (\lambda - d)$ steps are whole laps.}
\end{diagrambox}

\begin{alternativebox}[Approach 3: Binary Search on the Value (Counting)]
Forget the graph and apply the pigeonhole principle to a \emph{range of values}. Let $\text{cnt}(x)$ be the number of entries $\le x$ and $d$ the duplicate. Then $\text{cnt}(x) > x$ \textbf{iff} $d \le x$ (proof below), so binary search for the smallest $x$ with $\text{cnt}(x) > x$; each test is one read-only pass over the array.

\begin{lstlisting}[style=cpp]
#include <bits/stdc++.h>
using namespace std;

class Solution {
public:
    int findDuplicate(vector<int>& nums) {
        int n = nums.size() - 1, lo = 1, hi = n;   // the answer lies in [lo, hi]
        while (lo < hi) {
            int mid = (lo + hi) / 2, cnt = 0;
            for (int x : nums) if (x <= mid) cnt++; // read-only counting pass
            if (cnt > mid) hi = mid;                // too many small values: d <= mid
            else lo = mid + 1;                      // otherwise d > mid
        }
        return lo;
    }
};
\end{lstlisting}

\textbf{Why the test is monotone.} If $\text{cnt}(\textit{mid}) > \textit{mid}$, two of those entries share a value in $[1, \textit{mid}]$ (pigeonhole), so $d \le \textit{mid}$. If $\text{cnt}(\textit{mid}) \le \textit{mid}$, the other $n + 1 - \text{cnt} \ge n + 1 - \textit{mid}$ entries lie in $[\textit{mid}+1, n]$, which has only $n - \textit{mid}$ values: two of them collide, so $d > \textit{mid}$. Each step keeps $d$ inside $[lo, hi]$.

\begin{diagrambox}
\begin{tikzpicture}
    \dpindex{bx}{0.75}{$x$}{1}{9}
    \dparray{bc}{0}{$\text{cnt}(x)$}{1/cell, 2/cell, 3/cell, 4/cell, 5/curcell, 6/curcell, 8/anscell, 9/cell, 10/cell}
    \node[note, text width=3.6cm, align=left, anchor=west] at (8.4,0.35) {mid $= 5$: $5 \not> 5$, so $lo = 6$\\ mid $= 7$: $8 > 7$, so $hi = 7$\\ mid $= 6$: $6 \not> 6$, so $lo = 7$};
\end{tikzpicture}
\figcap{Counting on the trace input $[5,3,7,6,9,8,1,4,2,7]$. $\text{cnt}(x) > x$ holds exactly for $x \ge 7$; the three probed values are highlighted, and the search stops at $lo = hi = 7$.}
\end{diagrambox}

\begin{complexitybox}
\textbf{Time:} \bigO{n \log n} --- $\lceil \log_2 n \rceil \le 17$ passes of $n+1$ reads. \textbf{Space:} \bigO{1}, input untouched. \textbf{Verdict:} \ac{}. It satisfies all the rules too, but it is slower than Floyd and uses no graph structure at all; it is the solution to reach for if you do not spot the functional graph.
\end{complexitybox}
\end{alternativebox}

\begin{tricksbox}[Approach 4: Approaches That Break a Rule, and Brent's Variant]
\renewcommand{\arraystretch}{1.25}
\begin{tabularx}{\linewidth}{@{}>{\bfseries}l l X@{}}
Idea & Cost & Which rule it breaks \\
\midrule
\texttt{seen[]} / hash set & \bigO{n} time and memory & \bigO{n} extra memory (Algorithm~1 of the theory from $0$, or a plain ``seen before?'' scan) \\
sort, then neighbours & \bigO{n \log n} & sorting \texttt{nums} modifies it; sorting a copy needs \bigO{n} memory \\
negate \texttt{nums[|v|]} & \bigO{n}, \bigO{1} memory & writes into the input (the visited marks live in the signs) \\
\end{tabularx}

\medskip
\textbf{Brent's algorithm} (in words) is a second \bigO{1}-memory cycle finder. The hare walks one step at a time; the tortoise stays put and \emph{teleports} to the hare's position whenever the number of steps since the last teleport reaches a power of two ($1, 2, 4, \dots$). As soon as the hare lands on the tortoise, the number of steps since the last teleport is exactly $\lambda$. Then start two pointers at $0$, advance one of them $\lambda$ steps, and move both one step at a time: they meet at the entry after $\mu$ steps (the gap of $\lambda$ is one full lap). It usually evaluates $f$ fewer times than Floyd; on this problem it returns the same value (checked on the random tests of Section 6).
\end{tricksbox}

\subsubsection*{5. Trace}

Our own input with a longer tail and a separate cycle: $n = 9$, \;$\texttt{nums} = [5,3,7,6,9,8,1,4,2,7]$. The value $7$ is at indices $2$ and $9$; the walk from $0$ is $0,5,8,2,7,4,9,7,\dots$ with $\mu = 4$ and $\lambda = 3$. The program returns $7$.

\begin{diagrambox}
\begin{tikzpicture}
    \foreach \v/\x/\y in {0/0/0, 5/1.2/0, 8/2.4/0, 2/3.6/0, 7/4.8/0, 4/5.9/0.8, 9/5.9/-0.8, 1/8.4/0.8, 3/9.7/0, 6/8.4/-0.8}
        \node[gv] (g\v) at (\x,\y) {\v};
    \foreach \u/\w in {0/5, 5/8, 8/2, 2/7, 7/4, 4/9, 9/7, 1/3, 3/6, 6/1} \draw[gde] (g\u) -- (g\w);
\end{tikzpicture}
\figcap{The trace graph: tail $0 \to 5 \to 8 \to 2$, cycle $7 \to 4 \to 9 \to 7$ entered at $7$ (from $2$ and from $9$), and a separate cycle $1 \to 3 \to 6 \to 1$ that the walk from $0$ never reaches.}
\end{diagrambox}

\begin{tracebox}[{Phase 1 on $[5,3,7,6,9,8,1,4,2,7]$: T moves 1 step, H moves 2}]
\centering
\begin{tikzpicture}
    \begin{scope}
        \foreach \v/\x/\y/\s in {0/0/0/gvsrc, 5/1.1/0/gv, 8/2.2/0/gvcur, 2/3.3/0/gv, 7/4.4/0/gvseen, 4/5.3/0.75/gv, 9/5.3/-0.75/gv}
            \node[\s] (a\v) at (\x,\y) {\v};
        \foreach \u/\w in {0/5, 5/8, 8/2, 2/7, 7/4, 4/9, 9/7} \draw[gde] (a\u) -- (a\w);
        \node[gdist, above=1mm of a8] {T};
        \node[gdist, text=recursionred, above=1mm of a7] {H};
        \node[note, text width=5.6cm, align=left, anchor=west] at (6.2,0) {\textbf{Round 2:} T $= x_2 = 8$, H $= x_4 = 7$. The hare is already on the cycle, the tortoise still on the tail.};
    \end{scope}
    \begin{scope}[yshift=-2.8cm]
        \foreach \v/\x/\y/\s in {0/0/0/gvsrc, 5/1.1/0/gv, 8/2.2/0/gv, 2/3.3/0/gv, 7/4.4/0/gvcur, 4/5.3/0.75/gvseen, 9/5.3/-0.75/gv}
            \node[\s] (b\v) at (\x,\y) {\v};
        \foreach \u/\w in {0/5, 5/8, 8/2, 2/7, 7/4, 4/9, 9/7} \draw[gde] (b\u) -- (b\w);
        \node[gdist, above=1mm of b7] {T};
        \node[gdist, text=recursionred, right=1mm of b4] {H};
        \node[note, text width=5.6cm, align=left, anchor=west] at (6.2,0) {\textbf{Round 4:} T $= x_4 = 7$ reaches the entry ($\mu = 4$), H $= x_8 = 4$. Both are on the cycle now; the hare closes the gap by one step per round.};
    \end{scope}
    \begin{scope}[yshift=-5.6cm]
        \foreach \v/\x/\y/\s in {0/0/0/gvsrc, 5/1.1/0/gv, 8/2.2/0/gv, 2/3.3/0/gv, 7/4.4/0/gv, 4/5.3/0.75/gv, 9/5.3/-0.75/gvdst}
            \node[\s] (c\v) at (\x,\y) {\v};
        \foreach \u/\w in {0/5, 5/8, 8/2, 2/7, 7/4, 4/9, 9/7} \draw[gde] (c\u) -- (c\w);
        \node[gdist, text=accentorange, right=1mm of c9] {T = H};
        \node[note, text width=5.6cm, align=left, anchor=west] at (6.2,0) {\textbf{Round 6:} T $= x_6 = 9$, H $= x_{12} = 9$: they meet at $m = 9$. Indeed $t^\ast = 6$ is the smallest multiple of $\lambda = 3$ that is $\ge \mu = 4$.};
    \end{scope}
\end{tikzpicture}

\medskip
\small
\begin{tabular}{@{}l|cccccc@{}}
\toprule
round $t$ & 1 & 2 & 3 & 4 & 5 & 6 \\
\midrule
tortoise T $= x_t$ & 5 & 8 & 2 & 7 & 4 & \textbf{9} \\
hare H $= x_{2t}$ & 8 & 7 & 9 & 4 & 7 & \textbf{9} \\
\bottomrule
\end{tabular}
\end{tracebox}

\begin{tracebox}[{Phase 2: T restarts at 0, both move 1 step}]
\centering
\begin{tikzpicture}
    \begin{scope}
        \foreach \v/\x/\y/\s in {0/0/0/gvcur, 5/1.1/0/gv, 8/2.2/0/gv, 2/3.3/0/gv, 7/4.4/0/gv, 4/5.3/0.75/gv, 9/5.3/-0.75/gvseen}
            \node[\s] (d\v) at (\x,\y) {\v};
        \foreach \u/\w in {0/5, 5/8, 8/2, 2/7, 7/4, 4/9, 9/7} \draw[gde] (d\u) -- (d\w);
        \node[gdist, above=1mm of d0] {T};
        \node[gdist, text=recursionred, right=1mm of d9] {H};
        \node[note, text width=5.6cm, align=left, anchor=west] at (6.2,0) {\textbf{Step 0:} T $= 0$, H $= m = 9$. T is $\mu = 4$ steps from the entry; H is $\lambda - d = 1$ step from it, plus one lap of $3$.};
    \end{scope}
    \begin{scope}[yshift=-2.8cm]
        \foreach \v/\x/\y/\s in {0/0/0/gvsrc, 5/1.1/0/gv, 8/2.2/0/gvcur, 2/3.3/0/gv, 7/4.4/0/gv, 4/5.3/0.75/gvseen, 9/5.3/-0.75/gv}
            \node[\s] (e\v) at (\x,\y) {\v};
        \foreach \u/\w in {0/5, 5/8, 8/2, 2/7, 7/4, 4/9, 9/7} \draw[gde] (e\u) -- (e\w);
        \node[gdist, above=1mm of e8] {T};
        \node[gdist, text=recursionred, right=1mm of e4] {H};
        \node[note, text width=5.6cm, align=left, anchor=west] at (6.2,0) {\textbf{Step 2:} T $= 8$ (tail), H $= 4$ (cycle). A tail vertex never equals a cycle vertex, so no false meeting can happen before step $\mu$.};
    \end{scope}
    \begin{scope}[yshift=-5.6cm]
        \foreach \v/\x/\y/\s in {0/0/0/gvsrc, 5/1.1/0/gv, 8/2.2/0/gv, 2/3.3/0/gv, 7/4.4/0/gvdst, 4/5.3/0.75/gv, 9/5.3/-0.75/gv}
            \node[\s] (f\v) at (\x,\y) {\v};
        \foreach \u/\w in {0/5, 5/8, 8/2, 2/7, 7/4, 4/9, 9/7} \draw[gde] (f\u) -- (f\w);
        \node[gdist, text=accentorange, above=1mm of f7] {T = H};
        \node[note, text width=5.6cm, align=left, anchor=west] at (6.2,0) {\textbf{Step 4 $= \mu$:} both at $7$, the entry point. Return \textbf{7} --- indeed $\texttt{nums}[2] = \texttt{nums}[9] = 7$.};
    \end{scope}
\end{tikzpicture}

\medskip
\small
\begin{tabular}{@{}l|ccccc@{}}
\toprule
step $k$ & 0 & 1 & 2 & 3 & 4 \\
\midrule
T $= x_k$ & 0 & 5 & 8 & 2 & \textbf{7} \\
H $= x_{6+k}$ & 9 & 7 & 4 & 9 & \textbf{7} \\
\bottomrule
\end{tabular}
\end{tracebox}

The official examples run the same way. Example 1: phase 1 visits $(t,h) = (1,3), (3,4), (2,4), (4,4)$ and meets at $4$ after $t^\ast = 4$ rounds ($\mu = 3$, $\lambda = 2$); phase 2 gives $(0,4), (1,2), (3,4), (2,2)$: answer $2$. Example 2: $(3,4), (4,3), (2,2)$, then $(0,2), (3,3)$: answer $3$. Example 3: the first round already gives $(3,3)$, then $(0,3), (3,3)$: answer $3$. The counting search probes $\textit{mid} = 2$ ($\text{cnt} = 3$) and $\textit{mid} = 1$ ($\text{cnt} = 1$) in Example 1 and also returns $2$.

\subsubsection*{6. Test suite}

\begin{testbox}
\centering\small
\begin{tabular}{@{}lll@{}}
\toprule
\textbf{nums} & \textbf{Expected} & \textbf{Why this test} \\
\midrule
$[1,3,4,2,2]$ & 2 & official example 1: $\mu = 3$, 2-cycle \\
$[3,1,3,4,2]$ & 3 & official example 2: a separate self-loop at $1$ \\
$[3,3,3,3,3]$ & 3 & official example 3: value repeated 5 times \\
$[1,1]$ & 1 & smallest input ($n = 1$) \\
$[5,3,7,6,9,8,1,4,2,7]$ & 7 & the trace: $\mu = 4$, $\lambda = 3$, $t^\ast = 6$ \\
$a[i] = i+1$ ($i < n$), $a[n] = n$, $n = 10^5$ & 100000 & longest tail: $\mu = 10^5$, $\lambda = 1$ \\
$a[i] = i+1$ ($i < n$), $a[n] = 1$, $n = 10^5$ & 1 & longest cycle: $\mu = 1$, $\lambda = 10^5$ \\
\bottomrule
\end{tabular}
\end{testbox}

Floyd, the counting search and Brent's variant agree with the all-pairs brute force on the examples above and on $200\,000$ random arrays with $n \le 9$ (one value repeated $2..n+1$ times, the others distinct).

\subsubsection*{7. Common mistakes}

\begin{warningbox}[Mistake 1 --- starting the walk at a vertex that may lie on a cycle]
\begin{lstlisting}[style=snippet]
int t = nums[1], h = nums[nums[1]];  ...  t = 1;   // WRONG: vertex 1 may be on a cycle
int t = nums[0], h = nums[nums[0]];  ...  t = 0;   // RIGHT: nothing points to 0
\end{lstlisting}
Starting at $1$ returns $1$ for Example 2 (the self-loop) and $1$ for the trace input (cycle $1 \to 3 \to 6$) --- both wrong. Only vertex $0$ is guaranteed to have a tail, and only an entry point reached from a tail is a duplicate.
\end{warningbox}

\begin{warningbox}[Mistake 2 --- phase 1 and phase 2 from different starts]
\begin{lstlisting}[style=snippet]
int t = 0, h = 0;
do { t = nums[t]; h = nums[nums[h]]; } while (t != h);   // phase 1 from 0
t = nums[0];                                            // WRONG: phase 2 from nums[0]
\end{lstlisting}
The proof needs the same start in both phases. Here the restarted pointer is one step ahead, the two pointers circle the cycle at a constant offset and \textbf{never meet}: the loop runs forever on Examples 1 and 2 and on the trace input. Restart at $0$ --- or start \emph{both} phases at \texttt{nums[0]}, which also works: that is the $\rho$ of vertex \texttt{nums[0]}, with the same entry point.
\end{warningbox}

\begin{warningbox}[Mistake 3 --- \texttt{while} with both pointers at the start]
Writing \texttt{t = h = 0; while (t != h) \{\ldots\}} never enters phase 1, and phase 2 then returns $0$, which is not even a possible value. Either do one round before testing (as in Approach 2) or use \texttt{do \{\ldots\} while}.
\end{warningbox}

\begin{warningbox}[Mistake 4 --- binary search on positions instead of values]
\texttt{nums} is not sorted, so ``look at \texttt{nums[mid]}'' means nothing. Approach 3 searches the \emph{value range} $[1, n]$ and inspects the whole array at every step; the test is $\text{cnt}(\textit{mid}) > \textit{mid}$, with $\le$ inside the count.
\end{warningbox}

\subsubsection*{8. Takeaways}

\begin{itemize}[leftmargin=*]
    \item ``$n+1$ values in $[1, n]$'' $\Rightarrow$ the array is a functional graph on $0..n$ in which $0$ has no incoming arrow: the walk from $0$ always has a tail.
    \item The repeated value is the vertex with two incoming arrows on that walk --- the \textbf{entry point} of the $\rho$.
    \item Floyd: meet after $t^\ast$ rounds with $\lambda \mid t^\ast$; restart one pointer at the \emph{same} start, step both by one, and they meet after exactly $\mu$ steps at the entry. \bigO{n} time, \bigO{1} memory, read-only.
    \item When memory is no issue, the \texttt{first[]} array of Algorithm~1 is simpler; when the input is read-only and memory is \bigO{1}, pointers that walk are the only tool left.
\end{itemize}

\begin{variantsbox}
\variant{Entry of a cycle in a linked list}{\exmedium\ \lc{Medium}\ \varlink{Linked List Cycle II}{LeetCode 142}{https://leetcode.com/problems/linked-list-cycle-ii/}}
  {nodes with \texttt{next} pointers instead of an array; return the node where the cycle begins, or \texttt{null}.}
  {a functional graph with out-degree $\le 1$ (\texttt{null} = no arrow); the start is \texttt{head}, which may itself be on the cycle.}
  {the same two phases from \texttt{head}; phase 1 stops early when the hare reaches \texttt{null} (no cycle). The examples give the nodes at positions $1$, $0$ and \texttt{null}.}
\variant{Happy numbers}{\exeasy\ \lc{Easy}\ \varlink{Happy Number}{LeetCode 202}{https://leetcode.com/problems/happy-number/}}
  {no array at all: $x \mapsto$ sum of the squares of the digits of $x$; is $1$ eventually reached?}
  {an implicit functional graph on the integers; $1$ is a self-loop, so every walk ends in the cycle $\{1\}$ or in another cycle.}
  {Floyd on $f$ with \bigO{1} memory; the answer is ``meeting point $= 1$''. Examples: $19 \to$ true, $2 \to$ false; matches a hash-set walk for all $x \le 10^5$.}
\variant{Cycles in a circular array}{\exmedium\ \lc{Medium}\ \varlink{Circular Array Loop}{LeetCode 457}{https://leetcode.com/problems/circular-array-loop/}}
  {$i \to (i + \texttt{nums}[i]) \bmod n$; is there a cycle of length $> 1$ whose moves all go in the same direction?}
  {a functional graph; walks are cut as soon as the sign changes, so the usable graph has out-degree $\le 1$.}
  {Floyd from every start along same-sign arrows, rejecting self-loops; zero out each explored walk so the total stays \bigO{n}. Examples: true, false, true; matches a brute force on $10^5$ random arrays.}
\variant{Several values may repeat}{\varnone}
  {same array shape, but any number of values may repeat; return \emph{one} repeated value.}
  {unchanged: the walk from $0$ still has a tail, and its entry point still has two incoming arrows.}
  {Floyd still returns a repeated value (the entry point of $0$'s $\rho$); the counting search also still works, because its two pigeonhole arguments never use uniqueness. Both checked on $300\,000$ random arrays.}
\end{variantsbox}
